Competições que utilizam robôs móveis para percorrer e desenvolver labirintos já se encontram consolidadas há algum tempo. Tais competições têm grande importância no cenário internacional para estudo da robótica dado o seu caráter complexo e multidisciplinar. Os robôs móveis, \textit{micromouses}, são robôs com tração diferencial e de pequenas dimensões, daí o seu nome. Do projeto de concepção de um robô do tipo até a sua participação em competições existem diversas etapas que vão desde a escolha da sua geometria e dos seus componentes eletrônicos até ao desenvolvimento dos algorítmos de resolução de labirinto em si. Nesse contexto, o presente trabalho concentra-se na etapa de controle de velocidade, a qual é a etapa anterior ao seguimento de trajetórias e resolução de labirintos. Neste trabalho foi desenvolvida uma estratégia de comunicação sem fio, baseada no protocolo ESP-NOW, para que a coleta de informações comportasse a natureza móvel do robô e viabilizasse a telemetria em tempo real do robô. Em seguida foram projetadas malhas de controle -- utilizando o controlador PID -- internas e externas, de forma a orientar o comportamento individual das rodas e o comportamento global do robô. Em seguida os resultados foram analisados quanto às suas métricas de soma do erro, variância da saída e do sinal de controle, tempo de assentamento e sobressinal. Assim, é possível concluir que a adequada adoção de técnicas de medição e telemetria no projeto permite ao sistema pode ser controlado e acompanhado com sucesso e então é possível prosseguir para as etapas de aperfeiçoamento seguintes.

% Separe as palavras-chave por ponto
\palavraschave{Micromouse. Labirinto. Comunicação Sem Fio. Controle}